\documentclass[11pt, a4paper]{../problems_template}

\begin{document}

\begin{titlepage}
  \begin{center}
	\Huge{Московские математические регаты}
  \end{center}
\end{titlepage} 


\chapter{2016-2017}

\begin{exercise}
В правильном 21-угольнике 6 вершин покрашены красным цветом, а 7 вершин - синим. Обязательно ли найдутся два равных треугольника, один из которых с красными вершинами, а другой - с синими?
\end{exercise}

\end{document} 
